\section{Introduction}

Securitisation is the financial process of converting a pool of illiquid assets such as home loans, auto loans, education loans, or credit card receivables into tradable securities. Instead of waiting for borrowers to repay loans over several years, banks sell these loan pools to a Special Purpose Vehicle (SPV). The SPV issues securities to investors, who receive returns from the cash flows generated by borrowers.

This process improves liquidity in the banking system, transfers credit risk, and provides investment opportunities to institutional investors.

India follows a well-regulated securitisation framework governed primarily by the Reserve Bank of India (RBI) and the Securities and Exchange Board of India (SEBI). Rating agencies such as CRISIL, Moody's, and S\&P evaluate securitised assets based on credit quality, default probability, and expected losses.

The Day 1 objective of this project is to understand India's securitisation ecosystem, identify the regulatory framework, study rating methodologies, analyse the provided datasets, and map data fields required for future analytical models.